\chapter{Introduction}

Programming languages are the main interface between a programmer's ideas and a
machine's execution model. A language does not only describe what a program
should do; it also defines how programmers structure thought, how errors are
reported, how much control is exposed, how much work can be moved from run time
to compile time, and how the resulting program interacts with the operating
system. Compiler construction is therefore both a theoretical and practical
discipline. It combines formal language processing, type systems, intermediate
representations, data layout, code generation, run-time conventions, diagnostics,
testing, and engineering tradeoffs.

This Final Year Project is the design and implementation of \epl, a programming
language and compiler. The language is intentionally small, but it is not merely
a calculator or interpreter. It is a statically typed, expression-oriented,
single-file language that compiles ahead of time to native object code through
LLVM. It supports functions, local mutable variables, blocks, conditionals,
loops, structs, arrays, typed raw pointers, external function declarations,
integer casts, and compile-time evaluation. The compiler is implemented in Rust
and is organized as a sequence of explicit stages: lexing, parsing, typed tree
lowering, compile-time evaluation, lower IR generation, CFG cleanup, LLVM IR
emission, and native object generation.

The project was motivated by a common observation in modern language design:
many systems languages now include some form of compile-time computation. C++ has
\texttt{constexpr}, Zig has \texttt{comptime}, Rust has compile-time constants
and procedural macros, and many other languages provide macro systems or
specialized metaprogramming facilities. These mechanisms can remove repetitive
run-time work, generate data structures during compilation, improve type safety,
or help programmers express configuration in the language itself. However, a
compile-time execution mechanism also forces difficult questions. Which
operations are safe to execute during compilation? Can compile-time code call
ordinary functions? How are mutable local variables represented? How can the
compiler avoid accidentally depending on run-time state? How should the result be
lowered into executable code?

\epl{} answers these questions conservatively. A \texttt{comptime} expression is
first lowered and type checked as part of the normal program. It is then checked
for purity: it may use local compile-time variables, arithmetic, loops, structs,
arrays, and calls to functions marked \texttt{@pure}; it may not use string
literals, pointer operations, impure function calls, surrounding run-time
variables, or returns from the enclosing function. If the expression passes the
check, the compiler interprets it and replaces it with a constant. This keeps the
feature powerful enough to be useful while keeping the implementation small
enough to be studied in a final year project.

\section{Project Motivation}

The aim of the project is educational and experimental. A production language has
to solve a large number of secondary problems: package management, build
systems, separate compilation, standard libraries, target-specific ABIs,
tooling, long-term compatibility, editor support, optimization levels, and many
more. A final year compiler project has a different goal. It should make the
central ideas visible. The code should show how source text becomes tokens, how
tokens become syntax trees, how names and types are resolved, how structured
source constructs can be simplified into a smaller core language, how a
compile-time evaluator can share the same typed representation as the run-time
compiler, and how a lower representation can be mapped to LLVM.

\epl{} therefore follows a deliberately transparent architecture. The lexer is
hand-written instead of being generated by a tool. The parser is hand-written
recursive descent. The first intermediate representation, \irtree{}, remains
tree-shaped so that it is easy to inspect and compare with the source language.
The second intermediate representation is a control-flow graph with basic
blocks, definitions, instructions, terminators, and block arguments. LLVM
generation is explicit and uses the C API through the Rust \texttt{llvm-sys}
crate. The repository also includes documentation, source examples, expected
\irtree{} snapshots, and a small test harness.

The language is intentionally close to C in its low-level flavor. It exposes raw
pointers and external function declarations, and it does not implement ownership,
lifetimes, borrow checking, exceptions, garbage collection, modules, imports, or
generic types. At the same time, it is more expression-oriented than C: blocks,
\texttt{if}, and \texttt{loop} have result types, \texttt{return},
\texttt{break}, and \texttt{continue} are expressions of never type, and
compile-time evaluation is part of the typed expression system.

\section{Project Contribution}

The contribution of this thesis is the design and implementation of the \epl{}
language and compiler. The project contributes the following concrete artifacts:

\begin{itemize}
  \item a syntax and language reference for \epl{};
  \item a hand-written lexer that recognizes identifiers, keywords, literals,
        punctuation, whitespace, and line comments;
  \item a recursive descent parser that constructs an abstract syntax tree close
        to the source program;
  \item a type system covering unit, never, booleans, signed and unsigned
        integers, typed pointers, opaque pointers, arrays, and structs;
  \item semantic analysis that resolves names, checks function signatures,
        validates expressions, computes aggregate layouts, and lowers source
        constructs into a typed \irtree{};
  \item a purity checker and interpreter for compile-time expressions and
        \texttt{@pure} functions;
  \item a lower control-flow graph representation used for LLVM generation;
  \item cleanup passes that remove zero-sized operations and simplify basic
        control-flow patterns;
  \item LLVM IR generation and native object-file emission;
  \item a command-line interface for inspecting the AST, \irtree{}, lower IR,
        CFG, LLVM IR, and object output;
  \item a regression test harness using \irtree{} snapshots and lower IR smoke
        tests.
\end{itemize}

\begin{figure}[H]
  \centering
  \fbox{\begin{minipage}{0.88\textwidth}
  \centering
  \begin{tabular}{c}
  \texttt{source.epl} \\
  $\downarrow$ \\
  lexer: tokens \\
  $\downarrow$ \\
  parser: AST \\
  $\downarrow$ \\
  semantic analysis: typed \irtree{} \\
  $\downarrow$ \\
  compile-time evaluation and tree cleanup \\
  $\downarrow$ \\
  CFG lower IR \\
  $\downarrow$ \\
  LLVM IR \\
  $\downarrow$ \\
  native object file
  \end{tabular}
  \end{minipage}}
  \caption{High-level pipeline of the \epl{} compiler}
  \label{fig:pipeline}
\end{figure}

\section{Scope and Limitations}

The scope of the implementation is a working compiler for one source file at
a time. It is not a complete production ecosystem. The most important current
limitations are that there are no modules or imports, the C ABI is implemented
only for simple external calls, there are no pointer safety checks, there are no
array bounds checks, there is no name mangling or namespaces, and optimizations
are intentionally limited.

These limitations are part of the design space - they show which features
are essential to the central compiler pipeline. For example, a full C ABI
would require target-specific calling convention work, modules and separate
compilation would change the structure of name resolution. Advanced optimization
would require more development time. The project focuses instead on a complete,
coherent path from source to object code, as well as compile time evaluation.

\section{Thesis Organization}

The rest of this thesis is organized as follows. Chapter
\ref{chap:background} introduces compiler concepts used throughout the thesis.
Chapter \ref{chap:language} presents the \epl{} language design and semantics.
Chapter \ref{chap:architecture} describes the compiler architecture and command
line interface. Chapter \ref{chap:front-end} explains lexical analysis, parsing,
and AST construction. Chapter \ref{chap:semantic} documents semantic analysis,
type checking, \irtree{} lowering, and tree cleanup. Chapter
\ref{chap:comptime} focuses on compile-time evaluation. Chapter
\ref{chap:lower-ir} describes lower IR, CFG construction, and optimization
passes. Chapter \ref{chap:llvm} explains LLVM generation and linking. Chapter
\ref{chap:testing} presents examples, test results, and evaluation. Chapter
\ref{chap:conclusion} concludes the thesis and discusses future work. The
appendices include the grammar, selected examples, and representative generated
output.
